% AUTO-GENERATED by update_report_from_results.py -- do not edit by hand
\providecommand{\GeneratedForecastingInterpretation}{}
\renewcommand{\GeneratedForecastingInterpretation}{%
The deployable forecast-only model reaches 10.07~mg/dL MAE and 16.02~mg/dL RMSE on the held-out test split, with bias -1.16~mg/dL and nominal 80\% interval coverage of 79.2\%. The terminal 60-minute MAE is 14.52~mg/dL, which is larger than the multi-horizon average and should be reported separately. Compared with persistence on the exact same anchors, the model improves MAE by 1.64~mg/dL overall and by 2.50~mg/dL at 60 minutes.}
\providecommand{\GeneratedPersonalizationInterpretation}{}
\renewcommand{\GeneratedPersonalizationInterpretation}{%
The unmatched warm-up sweep shows lower MAE at longer warm-up lengths, but the evaluated anchor set shrinks with warm-up time. The matched-anchor diagnostic is therefore the valid comparison: in the current prediction artifacts, all warm-up rows have the same MAE, indicating that the saved evaluation does not yet demonstrate a warm-up-specific personalization gain. Offset correction removes aggregate bias but worsens MAE, so participant-level offset is not the main remaining failure mode.}
\providecommand{\GeneratedScenarioInterpretation}{}
\renewcommand{\GeneratedScenarioInterpretation}{%
Forecast-only, factual-future, meal-proxy, activity-proxy, and sleep/rest-proxy metrics are nearly identical. The largest mean absolute proxy-scenario delta is 0.0071~mg/dL, far below a clinically meaningful glucose change. The scenario branch is therefore currently inactive or weakly used, and proxy outputs should be treated only as diagnostics. AI-READI has no timed insulin dose or insulin-on-board inputs; \texttt{med\_insulin} is static metadata, not an editable action.}
\providecommand{\GeneratedSafetyInterpretation}{}
\renewcommand{\GeneratedSafetyInterpretation}{%
Median hypoglycemia alarms are conservative: precision is high but recall is low. The q0.1 risk rule substantially increases hypoglycemia recall at the cost of lower precision, which is the expected clinical safety trade-off. Hyperglycemia is more predictable than hypoglycemia, but the q0.9 risk rule similarly trades precision for sensitivity.}
\providecommand{\GeneratedSubgroupInterpretation}{}
\renewcommand{\GeneratedSubgroupInterpretation}{%
Participant-level subgroup metrics show a consistent phenotype gradient: insulin-dependent and higher-HbA1c participants are harder to forecast than healthy or lower-HbA1c participants. Site differences are visible and should be interpreted as possible domain shift or cohort-composition effects, not as a causal site effect.}
\providecommand{\GeneratedRuntimeInterpretation}{}
\renewcommand{\GeneratedRuntimeInterpretation}{%
The stateful Mamba run is feasible on a single GPU. The best validation checkpoint occurs at epoch 5; later epochs improve training loss but do not improve validation pinball, so the best checkpoint rather than the final epoch is used for evaluation. The ablation launcher prepares short probes, but full ablations remain a follow-up experiment.}
